\section{Conclusion and Open Problems}\label{sec:fut}
Theorem~\ref{thm:sf} and our previous work~\cite{gdmeasure,ppl2024} give a decidable characterisation of the \(\sf\)-measurable regular languages.
Its algebraic characterisation (Theorem~\ref{thm:algsf}) can be considered as a natural extension of Sch\"utzenberger's theorem for star-free languages.
The decidability of the \(\group\)-measurability for regular languages is still open.
Perhaps the profinite method can be useful to solve this problem.
The probabilistic independence of star-free and group languages describes how these two classes are orthogonal.
%It is interesting that whether there exists a local variety of regular languages \(\CC\) properly containing \(\sf\) and the probabilistic independence of \(\CC\) and \(\group\) also holds.
%If there is no such class, one can regard that the class of star-free languages is a ``maximal'' subclass of regular languages which is orthogonal to the class of group languages in the probabilistic sense.

In 2023, the first author posed a question whether the equation
\(\reg = \bool(\rext(\sf) \cup \rext(\group))\) holds or not, at some domestic workshops in Japan.
If this kind of equation \emph{were} true, it can be considered as a ``measure-theoretic decomposition theorem'' of regular languages.
Later, the second author claimed that a weaker equation
\(\reg = \rext(\bool(\sf \cup \group))\) holds and
this result was announced at the international conferences Highlights 2024 by the first author.
However, after some discussion between these two authors, a flaw was found in the proof of \(\reg = \rext(\bool(\sf \cup \group))\).
Actually, this equation does \emph{not} hold as follows.
\begin{proposition}\label{prop:counterexample}
Let \(\CT_4\) be the full transformation monoid on the four elements set \(\{1,2,3,4\}\) whose identity element is the identity function \(\mathrm{id}\).
Define \(e, f \in \CT_4\) as
\begin{align*}
 e\colon 1 \mapsto 1, 2 \mapsto 2, 3 \mapsto 1, 4 \mapsto 2 \quad \text{ and } \quad
 f\colon 1 \mapsto 4, 2 \mapsto 3, 3 \mapsto 3, 4 \mapsto 4.
\end{align*}
Let \(M\) be the submonoid of \(\CT_4\) generated by \(e\) and \(f\).
For an alphabet \(A = \{a,b\}\), define a morphism \(\eta\colon A^* \rightarrow M\) as \(\eta(a) = e\) and \(\eta(b) = f\).
Then \(\eta^{-1}(e)\) is \(\bool(\sf_A \cup \group_A)\)-immeasurable.
In particular, we have \(\reg_A \supsetneq \rext_A(\bool(\sf_A \cup \group_A))\).
\end{proposition}
The authors are still interested in whether there is a suitable notion \(\CO\colon 2^{A^*} \rightarrow 2^{A^*}\) of ``operation'' on languages, which is richer than the Boolean closure operator \(\bool\) but the simpler than the operation on languages corresponds to the wreath product operation on monoids, for a measure-theoretic decomposition theorem of regular languages like
 \(\reg = \CO(\rext(\sf) \cup \rext(\group))\) or
 \(\reg = \rext(\CO(\sf \cup \group))\).
